\documentclass[11pt]{article}
\usepackage[a4paper,margin=2.5cm]{geometry}
\usepackage{amsmath,amssymb}
\usepackage{booktabs}
\usepackage{longtable}
\usepackage[hidelinks]{hyperref}

\title{Appendix: From Tracer Measurements to DEB Parameters\\(Kaun et al., 2007)}
\author{}
\date{}

\begin{document}
\maketitle

This appendix links the glucose tracer assay in Kaun et al. (2007) to the DEB-consistent quantities used in our diet reconstruction and simulations. Because the paper reports key values graphically, the numeric entries below are digitized from Fig.~3B and are approximate.

\section*{A1. Experimental measurements (source data)}
Kaun et al. (2007) quantified glucose intake and post-ingestive retention in mid-third instar larvae using a radiolabeled tracer. Larvae fed for 15\,min on dead yeast paste (2:1 water:yeast, w/w) containing 0.08\% Brilliant Blue R dye and $2\,\mu\mathrm{Ci\,ml^{-1}}$ $^{14}$C-6-glucose (specific activity $58\,\mu\mathrm{Ci\,mmol^{-1}}$). After feeding, larvae were transferred to unlabeled medium for 3\,h to clear gut contents; retained $^{14}$C label was quantified by scintillation counting. The paper reports intake and retention in fmol of $^{14}$C-glucose per larva (Fig.~3B).

\section*{A2. Diet context and tracer role}
The tracer is a measurement tool, not an energetic substrate. Because the tracer mass/volume/density is not reported, its energetic contribution is treated as negligible. For DEB energetics we therefore use yeast paste as the sole energy-bearing substrate, with energy density taken from \texttt{diets\_kaun\_2007.tex} (Section~4: Diet 2: Yeast paste). The tracer measurements are used only to infer genotype-specific modifiers of effective assimilation; they do not set absolute food energy.

\section*{A3. Digitization and definitions}
Values in Fig.~3B were digitized from the published plots. We define the post-ingestive retention (absorption) efficiency as
\begin{equation}
  \eta_{\mathrm{abs}} = \frac{\text{Retention}}{\text{Intake}},
\end{equation}
where ``Retention'' is the $^{14}$C label remaining after the 3\,h clearing period.

\section*{A4. Digitized intake, retention, and efficiency}
\begin{center}
\begin{tabular}{llll}
\toprule
Genotype & Intake (fmol\,larva$^{-1}$) & Retention (fmol\,larva$^{-1}$) & $\eta_{\mathrm{abs}}$ \\
\midrule
for$^R$    & 600  & 309.68 & 0.52 \\
for$^S$    & 1200 & 193.55 & 0.16 \\
for$^{S2}$ & 1050 & 182.49 & 0.17 \\
\bottomrule
\end{tabular}
\end{center}

These values reproduce the qualitative result reported by Kaun et al. (2007): rovers retain roughly half of ingested tracer, whereas sitters retain about one-sixth.

\section*{A5. Conversion to C-moles and energy (tracer and yeast-proxy)}
For carbohydrates (glucose/fructose) we use the chemical potential
\begin{equation}
  \mu_{\mathrm{carb}} = 4.68\times10^{5}\,\mathrm{J\,C\text{-}mol^{-1}}.
\end{equation}
With $n_C=6$ and $1\,\mathrm{fmol}=10^{-15}\,\mathrm{mol}$, a value $B$ (fmol glucose per larva) converts to
\begin{align}
  n_{C,\mathrm{tr}} &= 6B\times10^{-15}\quad [\mathrm{C\text{-}mol}],\\
  E_{\mathrm{tr}} &= \mu_{\mathrm{carb}}\,n_{C,\mathrm{tr}}\quad [\mathrm{J}].
\end{align}

Using the retention values above:
\begin{center}
\begin{tabular}{llll}
\toprule
Genotype & Retention (fmol) & $n_{C,\mathrm{tr}}$ (C-mol\,larva$^{-1}$) & $E_{\mathrm{tr}}$ (J\,larva$^{-1}$) \\
\midrule
for$^R$    & 309.68 & $1.86\times10^{-12}$ & $8.70\times10^{-7}$ \\
for$^S$    & 193.55 & $1.16\times10^{-12}$ & $5.43\times10^{-7}$ \\
for$^{S2}$ & 182.49 & $1.09\times10^{-12}$ & $5.12\times10^{-7}$ \\
\bottomrule
\end{tabular}
\end{center}

These energies are tracer-linked and are negligible compared to the yeast-paste energy budget. To connect tracer retention to yeast energetics, we introduce a mapping factor
\begin{equation}
  \phi_Y \equiv \frac{n_{C,\mathrm{yeast\text{-}eq}}}{n_{C,\mathrm{tr}}},
\end{equation}
so that
\begin{equation}
  E_{Y,\mathrm{eq}} = \mu_Y\,\phi_Y\,n_{C,\mathrm{tr}}.
\end{equation}

If we adopt the working proxy $\phi_Y=1$ and $\mu_Y\approx\mu_{\mathrm{carb}}$, the yeast-equivalent retained carbon and energy are:
\begin{center}
\begin{tabular}{lll}
\toprule
Genotype & $n_{C,\mathrm{yeast\text{-}eq}}$ (C-mol\,larva$^{-1}$) & $E_{Y,\mathrm{eq}}$ (J\,larva$^{-1}$) \\
\midrule
for$^R$    & $1.86\times10^{-12}$ & $8.70\times10^{-7}$ \\
for$^S$    & $1.16\times10^{-12}$ & $5.43\times10^{-7}$ \\
for$^{S2}$ & $1.09\times10^{-12}$ & $5.12\times10^{-7}$ \\
\bottomrule
\end{tabular}
\end{center}

Otherwise, values scale by $(\mu_Y/\mu_{\mathrm{carb}})\phi_Y$.

\section*{A6. DEB-consistent interpretation}
We distinguish food density $X$ (C-mol\,cm$^{-3}$), chemical potential $\mu_X$ (J\,C-mol$^{-1}$), and the derived environmental energy density $\varepsilon_X=\mu_X X$ (J\,cm$^{-3}$). Only $X$ enters the functional response; $\varepsilon_X$ is a derived convenience when diets are specified by mass or molarity.

Assumptions used to map the tracer measurements into DEB parameters are:

\begin{enumerate}
  \item Retained tracer is treated as assimilated product in the measured pathway (post-ingestive retention $\rightarrow$ effective assimilation for carbohydrates).
  \item The genotype ranking from carbohydrate retention is used as a proxy for a unified effective assimilation efficiency applied to the yeast-paste diet.
  \item Tracer intake/retention is mapped to yeast intake/retention via a genotype-invariant proportionality factor $\phi_Y$ (defined above), because yeast-labeled measurements are not available.
\end{enumerate}

In DEB notation, the functional response must be computed explicitly,
\begin{equation}
  f = \frac{X}{X+K},
\end{equation}
so $f$ is not set a priori from ``high food'' arguments. The half-saturation coefficient can be written as
\begin{equation}
  K = \frac{\{\dot J_{Xm}\}}{\{\dot F_m\}} = \frac{\{\dot p_{Am}\}}{\kappa_X\,\mu_X\,\{\dot F_m\}},
\end{equation}
with $\{\dot p_{Am}\}=\kappa_X\,\mu_X\,\{\dot J_{Xm}\}$. If the maximum ingestion rate $\{\dot J_{Xm}\}$ (or $\{\dot p_{Xm}\}$) is identical across genotypes, then differences in effective assimilation efficiency $\kappa_X$ imply differences in $K$:
\begin{equation}
  \frac{K_i}{K_j} = \frac{\kappa_{X,j}}{\kappa_{X,i}}.
\end{equation}

Therefore the tracer data constrain relative genotype differences in effective assimilation (through $\eta_{\mathrm{abs}}$ and the mapping assumptions), while absolute energetic scaling is anchored by the yeast-paste diet reconstruction.

\end{document}
